\chapter{Configurations CLI FortiGate}
\label{annexe:cli}

\section{Configuration des interfaces}

\begin{lstlisting}[language=bash, caption={Interfaces FGT-Primary}]
config system interface
    edit "port1"
        set mode dhcp
        set allowaccess ping https ssh
    next
    edit "port2"
        set ip 192.168.10.1 255.255.255.0
        set allowaccess ping https ssh
    next
    edit "port3"
        set ip 169.254.0.1 255.255.255.252
        set allowaccess ping
    next
end
\end{lstlisting}

\section{Route par défaut}

\begin{lstlisting}[language=bash, caption={Route vers Internet}]
config router static
    edit 1
        set device "port1"
        set gateway 192.168.122.1
    next
end
\end{lstlisting}

\section{DHCP LAN1}

\begin{lstlisting}[language=bash, caption={Serveur DHCP FGT-Primary}]
config system dhcp server
    edit 1
        set interface "port2"
        set default-gateway 192.168.10.1
        set netmask 255.255.255.0
        config ip-range
            edit 1
                set start-ip 192.168.10.100
                set end-ip 192.168.10.200
            next
        end
        set dns-service default
    next
end
\end{lstlisting}

\section{Configuration HA Active-Passive}

\begin{lstlisting}[language=bash, caption={HA --- FGT-Primary (actif)}]
config system ha
    set group-name "FortiLab-HA"
    set mode a-p
    set hbdev "port3"
    set priority 200
    set session-sync-dev "port3"
end
\end{lstlisting}

\section{Politique Pare-feu}

\begin{lstlisting}[language=bash, caption={LAN1 vers WAN avec SNAT}]
config firewall policy
    edit 1
        set name "LAN1-to-WAN"
        set srcintf "port2"
        set dstintf "port1"
        set srcaddr "all"
        set dstaddr "all"
        set action accept
        set schedule "always"
        set service "HTTP" "HTTPS" "DNS" "PING"
        set logtraffic all
        set nat enable
    next
end
\end{lstlisting}
